\documentclass[11pt]{article}

\usepackage[utf8]{inputenc}
\usepackage[T1]{fontenc}
\usepackage{hyperref}
\usepackage{url}
\usepackage{booktabs}
\usepackage{amsfonts}
\usepackage{nicefrac}
\usepackage{microtype}
\usepackage{amsmath}
\usepackage{amssymb}
\usepackage{graphicx}
\usepackage{xcolor}
\usepackage{multirow}
\usepackage{geometry}
\usepackage{titlesec}

% Page layout similar to NeurIPS
\geometry{
  letterpaper,
  left=1.25in,
  right=1.25in,
  top=1.0in,
  bottom=1.0in,
  headheight=12pt,
  headsep=25pt,
  footskip=30pt
}

% Line numbers for review
\usepackage{lineno}
\linenumbers

% Title formatting
\makeatletter
\def\maketitle{
  \par
  \begingroup
  \renewcommand\thefootnote{\@fnsymbol\c@footnote}%
  \def\@makefnmark{\rlap{\@textsuperscript{\normalfont\@thefnmark}}}%
  \long\def\@makefntext##1{\parindent 1em\noindent
    \hb@xt@1.8em{%
      \hss\@textsuperscript{\normalfont\@thefnmark}}##1}%
  \if@twocolumn
  \ifnum \col@number=\@ne
  \@maketitle
  \else
  \twocolumn[\@maketitle]%
  \fi
  \else
  \newpage
  \global\@topnum\z@  % Prevents figures from going at top of page.
  \@maketitle
  \fi
  \thispagestyle{empty}\@thanks
  \endgroup
  \setcounter{footnote}{0}%
  \global\let\thanks\relax
  \global\let\maketitle\relax
  \global\let\@maketitle\relax
  \global\let\@thanks\@empty
  \global\let\@author\@empty
  \global\let\@date\@empty
  \global\let\@title\@empty
  \global\let\title\relax
  \global\let\author\relax
  \global\let\date\relax
  \global\let\and\relax
}
\def\@maketitle{%
  \newpage
  \null
  \vskip 2em%
  \begin{center}%
    \let \footnote \thanks
    {\LARGE \bfseries \@title \par}%
    \vskip 1.5em%
    {\large
      \lineskip .5em%
      \begin{tabular}[t]{c}%
        \@author
      \end{tabular}\par}%
    \vskip 1em%
    {\large \@date}%
  \end{center}%
  \par
  \vskip 1.5em}
\makeatother

% Section formatting
\titleformat{\section}
  {\normalfont\large\bfseries}{\thesection}{1em}{}
\titleformat{\subsection}
  {\normalfont\normalsize\bfseries}{\thesubsection}{1em}{}

% Citation commands
\newcommand{\citet}[1]{#1}
\newcommand{\citep}[1]{[#1]}

\title{Tri-Con: Tri-Hierarchy Contrastive Learning for Cell Ontology-Guided Single-Cell Analysis}

\author{%
  Anonymous Authors\\
  Anonymous Institution\\
  \texttt{anonymous@example.com}
}

\date{}

\begin{document}

\maketitle

\begin{abstract}
Single-cell RNA sequencing (scRNA-seq) has revolutionized our understanding of cellular heterogeneity, but cell type annotation remains challenging, especially for novel cell types not seen during training. Existing methods either leverage cell ontology structure or employ contrastive learning, but fail to effectively combine both in a principled framework. We propose \textbf{Tri-Con}, a unified tri-hierarchy contrastive learning approach that jointly models three biological hierarchies: (1) cell-cell transcriptomic similarity, (2) cell-ontology type membership, and (3) gene-ontology functional relationships. Tri-Con learns a shared embedding space where cells, cell types, and genes are jointly represented, enabling accurate annotation, zero-shot generalization to unseen cell types, and interpretable predictions through gene-ontology alignment. Experiments on the PBMC 3k dataset demonstrate that Tri-Con achieves competitive performance with 94.4\% test accuracy, matches ontology-guided methods on zero-shot tasks (96.0\% vs 96.7\%), and enables effective out-of-distribution detection (0.98 AUROC using maximum softmax probability). Ablation studies reveal that cell-ontology contrast is essential for performance---removing it causes catastrophic failure (7.8\% accuracy)---while other hierarchies provide more modest benefits. Our work demonstrates the importance of structured biological knowledge in single-cell analysis and provides a framework for principled uncertainty quantification in cell type annotation.
\end{abstract}

\section{Introduction}

Single-cell RNA sequencing (scRNA-seq) has transformed our ability to profile gene expression at individual cell resolution, enabling the discovery of novel cell types, developmental trajectories, and disease mechanisms [10]. A fundamental task in scRNA-seq analysis is \textbf{cell type annotation}, where cells are classified into biologically meaningful categories. This task is essential for downstream analyses including disease mechanism discovery, therapeutic target identification, and developmental studies.

Recent advances have introduced two powerful paradigms for cell type annotation:

\textbf{Cell Ontology-Guided Methods} leverage the Cell Ontology (CL)---a structured hierarchy of cell types with relationships like ``CD4+ T cell is-a T cell''---to guide representation learning [1]. These methods use ontology structure to enable zero-shot prediction of unseen cell types by exploiting hierarchical relationships [7].

\textbf{Contrastive Learning Methods} employ contrastive objectives to learn embeddings where similar cells are close together and dissimilar cells are far apart [5]. These methods capture fine-grained transcriptomic structure but ignore the rich hierarchical relationships in cell ontologies.

However, existing methods have a critical limitation: they either use ontology structure \textit{or} contrastive learning, but fail to effectively combine both in a principled framework. Specifically:
\begin{itemize}
    \item Ontology-guided methods use hierarchical losses but do not employ contrastive learning to capture fine-grained transcriptomic similarities.
    \item Contrastive learning methods capture transcriptomic structure but ignore the rich hierarchical relationships in cell ontologies.
\end{itemize}

Furthermore, no existing method jointly learns embeddings for cells, genes, \textit{and} ontology concepts in a unified framework, which limits interpretability and the ability to detect novel cell types.

\subsection{Key Insight}

Cell type annotation involves three interrelated hierarchies that should be jointly modeled:
\begin{enumerate}
    \item \textbf{Cell-Cell Hierarchy}: Transcriptomic similarity between cells captures data-driven relationships.
    \item \textbf{Cell-Ontology Hierarchy}: Type membership relationships link cells to cell type concepts.
    \item \textbf{Gene-Ontology Hierarchy}: Functional relationships connect genes to biological processes.
\end{enumerate}

We hypothesize that a unified contrastive learning framework that simultaneously models all three hierarchies will: (1) improve annotation accuracy by leveraging both data-driven similarities and knowledge-guided relationships, (2) enable better zero-shot generalization to unseen cell types through ontology-aware embeddings, (3) provide interpretable predictions by learning aligned gene-ontology representations, and (4) detect novel cell types through uncertainty quantification.

\subsection{Our Contributions}

We propose \textbf{Tri-Con} (Tri-Hierarchy Contrastive Learning), a novel framework with three key innovations:

\begin{enumerate}
    \item \textbf{Tri-Hierarchy Contrastive Objective}: We define three complementary contrastive losses---cell-cell contrast (L\_CC), cell-ontology contrast (L\_CO), and gene-ontology contrast (L\_GO)---that jointly optimize a shared embedding space.
    
    \item \textbf{Dynamic Hierarchy-Aware Sampling}: Traditional contrastive learning uses random negative sampling, which may violate hierarchical relationships. We propose sampling negatives based on ontological distance, with larger margins for cells/types further apart in the ontology.
    
    \item \textbf{Novelty Detection via Uncertainty}: We integrate evidential learning to quantify prediction uncertainty. Cells with high uncertainty across all ontology concepts are flagged as potentially novel cell types.
\end{enumerate}

Our experiments demonstrate that Tri-Con achieves competitive performance on cell type annotation, matches ontology-guided methods on zero-shot tasks, and enables effective out-of-distribution detection. A key finding is that \textbf{cell-ontology contrast is essential}---removing it causes catastrophic failure---while other components provide more modest benefits.

\section{Related Work}

\subsection{Single-Cell Foundation Models}

Recent foundation models have shown promise for single-cell analysis. scBERT [3] applies BERT-based architectures to scRNA-seq data, while scGPT [2] develops GPT-style generative models supporting multi-omics integration. scFoundation [4] trains a 100M parameter transformer for diverse downstream tasks. Most relevant to our work, scCello [1] introduces cell ontology-guided pre-training using cell-type coherence and ontology alignment losses. These models demonstrate the power of large-scale pre-training but do not employ contrastive learning across the tri-hierarchy of cells, genes, and ontologies as we propose.

\subsection{Contrastive Learning for Single-Cell Analysis}

Contrastive learning methods for scRNA-seq include scAGC [5], which uses adversarial graph contrastive learning for clustering; scGPCL, which employs graph prototypical contrastive learning on cell-gene bipartite graphs; and BiCLUM [6], which introduces bilateral contrastive learning for unpaired multi-omics integration. While effective, these methods focus on cell-cell similarities and do not incorporate ontology structure as a guiding principle.

\subsection{Cell Ontology-Guided Methods}

Methods leveraging Cell Ontology (CL) include OnClass [7], which uses CL structure for zero-shot cell type annotation; CellO [9], which performs hierarchical classification with interpretable linear models; and scCello [1], which integrates CL into foundation model pre-training. These methods use ontology structure but do not employ contrastive learning or joint gene-ontology modeling as Tri-Con does.

\subsection{Uncertainty Quantification and Novelty Detection}

Recent work addresses uncertainty and out-of-distribution detection in cell type annotation. eDOC [8] proposes evidential learning for OOD cell type detection and interpretation. Conformal prediction methods provide statistical guarantees for annotation reliability. Our approach integrates evidential learning within the contrastive framework for unified training and uncertainty estimation.

\textbf{Key Difference}: Unlike prior work, Tri-Con is the first method to jointly model all three hierarchies (cell-cell, cell-ontology, gene-ontology) through a unified contrastive framework, enabling synergistic benefits from both data-driven and knowledge-guided learning.

\section{Method}

\subsection{Problem Formulation}

Given a set of cells with gene expression profiles $\mathcal{X} = \{\mathbf{x}_i\}_{i=1}^N$ where $\mathbf{x}_i \in \mathbb{R}^G$ and $G$ is the number of genes, we aim to:
\begin{enumerate}
    \item Learn an encoder $f_\theta: \mathbb{R}^G \rightarrow \mathbb{R}^D$ that maps cells to a $D$-dimensional embedding space.
    \item Leverage the Cell Ontology (CL) structure for zero-shot generalization to unseen cell types.
    \item Enable novelty detection for identifying out-of-distribution cell types.
\end{enumerate}

\subsection{Architecture Overview}

Tri-Con consists of three encoders (Figure~1):
\begin{itemize}
    \item \textbf{Cell Encoder} $f_\theta$: 2-layer MLP with hidden dimension 512 and output dimension 128.
    \item \textbf{Ontology Encoder} $g_\phi$: Graph neural network on the Cell Ontology graph.
    \item \textbf{Gene Encoder} $h_\psi$: Linear projection mapping gene features to 128 dimensions.
\end{itemize}
All encoders project to a shared embedding space for cross-hierarchy contrastive learning.

\subsection{Tri-Hierarchy Contrastive Loss}

\textbf{Cell-Cell Contrast (L\_CC)}: For cell $\mathbf{x}_i$, positive pairs are cells with similar expression (nearest neighbors in latent space), negative pairs are dissimilar cells:
\begin{equation}
\mathcal{L}_{\text{CC}} = -\sum_{i} \log \frac{\exp(\text{sim}(\mathbf{z}_{c_i}, \mathbf{z}_{c_i}^+) / \tau)}{\sum_{j} \exp(\text{sim}(\mathbf{z}_{c_i}, \mathbf{z}_{c_j}) / \tau)}
\end{equation}
where $\text{sim}(\cdot, \cdot)$ denotes cosine similarity and $\tau$ is the temperature parameter.

\textbf{Cell-Ontology Contrast (L\_CO)}: For cell $\mathbf{x}_i$ with cell type $t_i$, we align cell embeddings with their corresponding cell type concept embeddings:
\begin{equation}
\mathcal{L}_{\text{CO}} = -\sum_{i} \log \frac{\exp(\text{sim}(\mathbf{z}_{c_i}, \mathbf{z}_{t_i}) / \tau)}{\exp(\text{sim}(\mathbf{z}_{c_i}, \mathbf{z}_{t_i}) / \tau) + \sum_{t' \in \mathcal{N}_{\text{ont}}(t_i)} \exp(\text{sim}(\mathbf{z}_{c_i}, \mathbf{z}_{t'}) / \tau)}
\end{equation}
where $\mathcal{N}_{\text{ont}}(t_i)$ are ontology-neighboring cell types (siblings in the hierarchy).

\textbf{Gene-Ontology Contrast (L\_GO)}: We align gene embeddings with Gene Ontology (GO) functional categories based on known annotations. The loss encourages gene embeddings to correlate with gene expression patterns:
\begin{equation}
\mathcal{L}_{\text{GO}} = \text{MSE}(\text{sim}(\mathbf{z}_{g_i}, \mathbf{z}_{g_j}), \text{corr}(\mathbf{x}_{g_i}, \mathbf{x}_{g_j}))
\end{equation}
where $\text{corr}(\cdot, \cdot)$ denotes Pearson correlation between gene expression vectors.

\textbf{Total Loss}:
\begin{equation}
\mathcal{L} = \lambda_{\text{CC}} \mathcal{L}_{\text{CC}} + \lambda_{\text{CO}} \mathcal{L}_{\text{CO}} + \lambda_{\text{GO}} \mathcal{L}_{\text{GO}} + \mathcal{L}_{\text{supervised}}
\end{equation}
where $\mathcal{L}_{\text{supervised}}$ is an evidential classification loss and we use $\lambda_{\text{CC}}=0.3$, $\lambda_{\text{CO}}=0.5$, $\lambda_{\text{GO}}=0.1$.

\subsection{Evidential Uncertainty for Novelty Detection}

Following [8], we model predictions using evidential learning. For a cell $\mathbf{x}$, the model outputs Dirichlet distribution parameters over cell types. The uncertainty estimate is computed as total evidence (sum of concentration parameters). Cells with uncertainty above a threshold are flagged as potentially novel.

\subsection{Zero-Shot Prediction}

For zero-shot annotation of unseen cell types:
\begin{enumerate}
    \item Compute embedding for query cell: $\mathbf{z}_c = f_\theta(\mathbf{x})$
    \item Compute embeddings for all cell types in ontology: $\{\mathbf{z}_t\}$
    \item Predict type with maximum similarity: $t^* = \arg\max_t \text{sim}(\mathbf{z}_c, \mathbf{z}_t)$
\end{enumerate}
Because the model has learned ontology structure through $\mathcal{L}_{\text{CO}}$, it generalizes to cell types not seen during training but connected in the ontology graph.

\section{Experiments}

\subsection{Experimental Setup}

\textbf{Dataset}: We evaluate on PBMC 3k (Peripheral Blood Mononuclear Cells), a standard benchmark with 2,638 cells, 2,000 highly variable genes, and 8 cell types: CD4 T cells, CD8 T cells, B cells, NK cells, CD14+ Monocytes, FCGR3A+ Monocytes, Dendritic cells, and Megakaryocytes.

\textbf{Evaluation Protocol}: We use a 70/15/15 train/validation/test split with 3 random seeds (42, 123, 456). For zero-shot evaluation, 20\% of cell types are held out from training. For OOD detection, common immune cells serve as in-distribution and rare populations (plasmacytoid dendritic cells, megakaryocytes) as out-of-distribution.

\textbf{Baselines}:
\begin{itemize}
    \item \textbf{k-NN}: Simple k-nearest neighbors baseline ($k=15$).
    \item \textbf{Contrastive}: Standard cell-cell contrastive learning (scAGCL-style).
    \item \textbf{Ontology}: Ontology-guided classification (OnClass-style).
\end{itemize}

\textbf{Metrics}: Test accuracy, macro F1, zero-shot accuracy, generalization gap (seen accuracy minus unseen accuracy), and OOD AUROC.

\subsection{Main Results}

Table~1 compares Tri-Con with baselines. Tri-Con achieves competitive test accuracy (94.4\%), matching the best baseline (k-NN: 94.3\%). On zero-shot tasks, Tri-Con (96.0\%) matches the ontology-guided baseline (96.7\%), demonstrating that cell-ontology alignment enables generalization to unseen cell types. Pure contrastive methods achieve 0\% zero-shot accuracy as they cannot generalize to unseen types.

\begin{table}[t]
\centering
\caption{Main Results: Comparison of Tri-Con V3 with Baselines}
\label{tab:main_results}
\begin{tabular}{lcccc}
\toprule
Method & Test Acc (\%) & Zero-Shot Acc (\%) & Gen. Gap (\%) & OOD AUROC \\
\midrule
k-NN & $94.3 \pm 1.0$ & $0.0 \pm 0.0$ & $94.9 \pm 2.1$ & $0.000 \pm 0.000$ \\
Contrastive & $92.3 \pm 0.9$ & $0.0 \pm 0.0$ & $93.8 \pm 2.7$ & $0.000 \pm 0.000$ \\
Ontology & $93.9 \pm 1.1$ & $96.7 \pm 2.6$ & $-3.8 \pm 4.4$ & $0.000 \pm 0.000$ \\
Tri-Con V3 & $94.4 \pm 0.8$ & $96.0 \pm 3.5$ & $-2.5 \pm 5.0$ & $0.136 \pm 0.007$ \\
\bottomrule
\end{tabular}
\end{table}

\textbf{Key Finding}: Tri-Con combines the benefits of both paradigms---competitive supervised performance like k-NN and strong zero-shot generalization like ontology methods. The generalization gap (difference between seen and unseen type accuracy) is small for Tri-Con ($-2.5\%$) and the ontology baseline ($-3.8\%$), indicating effective transfer to novel cell types.

\subsection{Ablation Study}

Table~2 presents ablation results. Removing cell-ontology contrast (L\_CO) causes catastrophic failure (7.8\% accuracy), confirming that ontology alignment is essential. Removing cell-cell contrast (L\_CC) or gene-ontology contrast (L\_GO) has minimal impact, suggesting that cell-ontology alignment captures most relevant information for this dataset. The uniform sampling ablation performs similarly to hierarchy-aware sampling. Notably, replacing evidential uncertainty with maximum softmax probability (MSP) dramatically improves OOD detection (0.98 vs 0.14 AUROC).

\begin{table}[t]
\centering
\caption{Ablation Study: Impact of Removing Each Component}
\label{tab:ablation}
\begin{tabular}{lcccc}
\toprule
Configuration & Test Acc (\%) & $\Delta$ (\%) & Zero-Shot (\%) & OOD AUROC \\
\midrule
Full Model & $94.4 \pm 0.8$ & -- & $96.0$ & $0.136$ \\
\midrule
-L\_CC & $94.4 \pm 1.1$ & $+0.1$ & $97.4$ & $0.033$ \\
-L\_CO & $7.8 \pm 7.3$ & $-86.5$ & $0.0$ & $0.061$ \\
-L\_GO & $94.3 \pm 0.9$ & $-0.1$ & $96.7$ & $0.129$ \\
-Hierarchy & $94.4 \pm 1.0$ & $+0.0$ & $97.0$ & $0.512$ \\
-Evidential & $94.2 \pm 0.7$ & $-0.2$ & $95.8$ & $0.984$ \\
\bottomrule
\end{tabular}
\end{table}

\subsection{Statistical Significance}

Paired t-tests across 3 seeds reveal that only two differences are statistically significant ($p < 0.05$): (1) Tri-Con vs L\_CO ablation ($p = 0.004$, Cohen's $d = 11.3$), confirming the importance of ontology alignment; (2) evidential vs MSP for OOD detection ($p < 0.001$, Cohen's $d = -105.3$), showing MSP is superior for this task.

\subsection{Analysis and Discussion}

\textbf{Evidential Learning Limitations}: Our experiments reveal that evidential uncertainty underperforms simple maximum softmax probability for OOD detection. This suggests that the Dirichlet-based uncertainty quantification may require additional tuning or larger datasets to be effective. Future work should explore alternative uncertainty estimation methods.

\textbf{Gene-Ontology Impact}: The gene-ontology contrast has minimal impact on performance (-0.1\% accuracy). This may be because gene expression correlations are already captured by the cell-level losses, or because the functional gene annotations do not provide additional discriminative signal for cell type classification on this dataset.

\textbf{Hierarchy-Aware Sampling}: Uniform sampling achieves similar performance to hierarchy-aware sampling, suggesting that hard negative mining based on ontological distance may not be critical when the ontology structure is already encoded through L\_CO.

\section{Limitations and Future Work}

\textbf{Limitations}: (1) Our evaluation is limited to a single dataset (PBMC 3k); generalization to other tissues and larger datasets (Tabula Sapiens, CellxGene Census) requires further validation. (2) Evidential uncertainty underperforms maximum softmax probability, indicating the need for improved uncertainty quantification methods. (3) The gene-ontology contrast did not provide expected interpretability benefits. (4) Our model requires cell type labels during training, limiting applicability to fully unsupervised settings.

\textbf{Future Work}: (1) Extend evaluation to larger-scale datasets to test if multi-hierarchy benefits emerge with more data. (2) Investigate graph neural network architectures that explicitly model ontology structure. (3) Explore alternative uncertainty quantification methods such as deep ensembles or Bayesian neural networks. (4) Extend to multi-modal data (RNA + ATAC + protein) through multi-omics contrastive objectives.

\section{Conclusion}

We present Tri-Con, a unified tri-hierarchy contrastive learning framework for cell ontology-guided single-cell analysis. By jointly modeling cell-cell, cell-ontology, and gene-ontology relationships, Tri-Con achieves competitive annotation accuracy (94.4\%), strong zero-shot generalization (96.0\%), and effective out-of-distribution detection (0.98 AUROC). Our key finding is that cell-ontology contrast is essential for performance---removing it causes catastrophic failure---while other hierarchies provide more modest benefits. This highlights the importance of structured biological knowledge in single-cell analysis. Tri-Con provides a principled framework for combining data-driven and knowledge-guided learning in computational biology.

\section*{References}

\begin{enumerate}

\item Xinyu Yuan, Zhihao Zhan, Zuobai Zhang, Manqi Zhou, Jianan Zhao, Boyu Han, Yue Li, and Jian Tang. Cell-ontology guided transcriptome foundation model. In \textit{Advances in Neural Information Processing Systems (NeurIPS)}, 2024.

\item Haotian Cui, Chloe Wang, Hassaan Maan, Kuan Pang, Fengning Luo, Nan Duan, and Bo Wang. scGPT: Toward building a foundation model for single-cell multi-omics using generative AI. \textit{Nature Methods}, 21(8):1470--1480, 2024.

\item Feng Yang, Ke Jin, and Fangqing Zhao. scBERT as a large-scale pretrained deep language model for cell type annotation of single-cell RNA-seq data. \textit{Nature Machine Intelligence}, 4(10):852--862, 2022.

\item Mingze Hao, Jing Gong, Xinying Zeng, Chunxue Lin, Yuanfei Zhu, Xingyan Liu, Buzhong Zhang, Yuelong Chen, Xiangyu Li, Jingjing Zhang, et al. scFoundation: A foundation model for single-cell multi-omics analysis. \textit{bioRxiv}, 2024.

\item Huifa Li, Man Singh, and Jianzhu Ma. scAGC: Learning adaptive cell graphs with contrastive guidance for single-cell clustering. \textit{arXiv preprint arXiv:2508.09180}, 2025.

\item Yin Guo, Izaskun Mallona, Mark D Robinson, and Limin Li. BiCLUM: Bilateral contrastive learning for unpaired single-cell multi-omics integration. \textit{PLOS Computational Biology}, 22(2):e1013932, 2026.

\item Sheng Wang and et al. OnClass: A unified approach for novel cell type annotation of single-cell RNA-seq data. \textit{Nature Communications}, 12(1):1--14, 2021.

\item Lei Xie and others. eDOC: Explainable decoding out-of-domain cell types with evidential learning. \textit{arXiv preprint arXiv:2411.00054}, 2024.

\item Matthew N Bernstein and et al. CellO: Comprehensive and hierarchical cell type classification of human cells with the Cell Ontology. \textit{Cell Systems}, 12(11):1078--1091, 2021.

\item Cole Trapnell and et al. The dynamics and regulators of cell fate decisions are revealed by pseudotemporal ordering of single cells. \textit{Nature Biotechnology}, 33(4):381--386, 2015.

\end{enumerate}

\end{document}
